Taking account of the definition of @@M0@@, and of (6.6.4), we obtain:
@@M1@@
By (6.5.3), and the fact that @@M2@@, we finally obtain:
@@M3@@
whence the result.

\sgaStableTarget{sga6.e05.prop.6-10}\textbf{Proposition 6.10.} \uline{Let} @@M4@@ \uline{be a binomial ring}, @@M5@@ \uline{an augmented} @@M6@@-@@M7@@-\uline{algebra}, @@M8@@ \uline{an element of} @@M9@@ \uline{such that} @@M10@@ \uline{for} @@M11@@, \uline{and} @@M12@@. \uline{Then}:
@@M13@@

We can find elements @@M14@@, and @@M15@@ such that:
@@M16@@
with @@M17@@, and @@M18@@ for every @@M19@@ and every @@M20@@. Let @@M21@@ be the @@M22@@-@@M23@@-algebra generated by generators @@M24@@, @@M25@@ subject to the relations:
@@M26@@
We define an augmentation homomorphism @@M27@@ by @@M28@@. By (4.9.1), @@M29@@ is the ring of polynomials with coefficients in @@M30@@ in the indeterminates @@M31@@, with @@M32@@, and @@M33@@ with @@M34@@, and the @@M35@@-filtration of @@M36@@ is identical to its weight filtration, assigning the weight @@M37@@ to @@M38@@ and @@M39@@ (4.12 and 4.15).

We put:
@@M40@@
